\section{切向加速度/法向加速度}

在自然坐标系中，速度 $\vec{v}$ 的方向永远沿着轨道切线方向。速度矢量可以写为大小（速率 $v$）与方向（切向单位矢量 $\vec{\tau}$）的乘积：$\vec{v} = v\vec{\tau}$，其中 $\vec{\tau}$ 是关于时间 $t$ 的函数（由于质点曲线运动时，切向单位矢量 $\vec{\tau}$ 的方向一直在改变）。对速度矢量应用乘积求导法则： \(\vec{a} = \frac{\mathrm{d}\vec{v}}{\mathrm{d}t} = \frac{\mathrm{d}(v\vec{\tau})}{\mathrm{d}t} = \frac{\mathrm{d}v}{\mathrm{d}t}\vec{\tau} + v\frac{\mathrm{d}\vec{\tau}}{\mathrm{d}t}\) 
这一步非常直观地把加速度分成了两部分：
\begin{itemize}
\item 第一项 $\frac{\mathrm{d}v}{\mathrm{d}t}\vec{\tau}$：速度大小随时间的变化率，方向沿切向。
\item 第二项 $v\frac{\mathrm{d}\vec{\tau}}{\mathrm{d}t}$：速度方向随时间的变化率；还需要确定 $\frac{\mathrm{d}\vec{\tau}}{\mathrm{d}t}$。
\end{itemize}
考察 $\mathrm{d}\vec{\tau}$ 与路程微元 $\mathrm{d}s$ 的关系，链式展开后分离出速率 $v$：
 \(\frac{\mathrm{d}\vec{\tau}}{\mathrm{d}t} = \frac{\mathrm{d}\vec{\tau}}{\mathrm{d}s} \cdot \frac{\mathrm{d}s}{\mathrm{d}t} = \frac{\mathrm{d}\vec{\tau}}{\mathrm{d}s} v\) 
设质点在 $\mathrm{d}t$ 内走过弧长 $\mathrm{d}s$，对应的切向单位矢量由 $\vec{\tau}$ 变为 $\vec{\tau}'$。由于它们都是单位矢量，模长保持不变：$|\vec{\tau}| = |\vec{\tau}'| = 1$，将 $\vec{\tau}$ 和 $\vec{\tau}'$ 移到同一起点，它们之间的夹角为 $\mathrm{d}\theta$（即轨道切线转过的角度）。
\begin{itemize}
\item 方向：当 $\mathrm{d}\theta \to 0$ 时，矢量差 $\mathrm{d}\vec{\tau} = \vec{\tau}' - \vec{\tau}$ 将垂直于 $\vec{\tau}$，严格指向曲线内凹的一侧，即法向单位矢量 $\vec{n}$ 的方向。
\item 大小：由大圆弧或三角形小角近似关系，当 $\mathrm{d}\theta$ 极小时，两单位矢量终点连线的长度（即 $|\mathrm{d}\vec{\tau}|$）等于弧长： \(|\mathrm{d}\vec{\tau}| = |\vec{\tau}| \cdot \mathrm{d}\theta = 1 \cdot \mathrm{d}\theta = \mathrm{d}\theta\) 
\end{itemize}
根据轨道曲率半径 $\rho$ 的定义，弧长 $\mathrm{d}s$、转角 $\mathrm{d}\theta$ 与半径 $\rho$ 满足：
 \(\mathrm{d}s = \rho \mathrm{d}\theta \implies \mathrm{d}\theta = \frac{\mathrm{d}s}{\rho}\) 
代入得： \(|\mathrm{d}\vec{\tau}| = \frac{\mathrm{d}s}{\rho}\) 
综合上面的方向（$\vec{n}$）和大小（$\frac{\mathrm{d}s}{\rho}$），可以写出 $\mathrm{d}\vec{\tau}$ 的严格矢量表达式： \(\mathrm{d}\vec{\tau} = \frac{\mathrm{d}s}{\rho}\vec{n} \implies \frac{\mathrm{d}\vec{\tau}}{\mathrm{d}s} = \frac{1}{\rho}\vec{n}\) 
将 $\frac{\mathrm{d}\vec{\tau}}{\mathrm{d}s} = \frac{1}{\rho}\vec{n}$ 代回链式法则： \(\frac{\mathrm{d}\vec{\tau}}{\mathrm{d}t} = \frac{\mathrm{d}\vec{\tau}}{\mathrm{d}s} v = \left(\frac{1}{\rho}\vec{n}\right) v = \frac{v}{\rho}\vec{n}\) 再代入加速度展开式：
\begin{definition}[切向加速度/法向加速度]
$$\vec{a} = \frac{\mathrm{d}\vec{v}}{\mathrm{d}t} = \frac{\mathrm{d}(v\vec{\tau})}{\mathrm{d}t} = \frac{\mathrm{d}v}{\mathrm{d}t}\vec{\tau} + v\frac{\mathrm{d}\vec{\tau}}{\mathrm{d}t} = \frac{\mathrm{d}v}{\mathrm{d}t}\vec{\tau} + v\left(\frac{v}{\rho}\vec{n}\right)= \frac{\mathrm{d}v}{\mathrm{d}t}\vec{\tau} + \frac{v^2}{\rho}\vec{n}$$
\begin{itemize}
\item $\vec{a}$: \textbf{加速度矢量（Total Acceleration Vector）}。质点在空间运动时的总瞬时加速度。
\item $\vec{v}$: \textbf{速度矢量（Velocity Vector）}。其方向永远指向质点运动轨迹的切线方向。
\item $v$: \textbf{瞬时速率（Speed）}。是一个标量，代表速度矢量 $\vec{v}$ 的模长（即 $v = |\vec{v}| = \frac{\mathrm{d}s}{\mathrm{d}t}$），反映质点运动的快慢。
\item $\vec{\tau}$: \textbf{切向单位矢量（Unit Tangent Vector）}。是一个沿轨道切线、指向质点前进方向的单位矢量（满足 $|\vec{\tau}| = 1$），负责指示速度的方向。
\item $\vec{n}$: \textbf{法向单位矢量（Unit Normal Vector）}。是一个垂直于切线、严格指向轨道内凹一侧（曲率中心）的单位矢量（满足 $|\vec{n}| = 1$）。
\item $\rho$: \textbf{曲率半径（Radius of Curvature）}。轨迹在该点处局部近似圆弧的半径。轨迹越弯曲，$\rho$ 越小；直线运动时 $\rho \to \infty$。
\item $\frac{\mathrm{d}v}{\mathrm{d}t}\vec{\tau}$: \textbf{切向加速度矢量（Tangential Acceleration）}。通常记作 $\vec{a}_\tau$。其大小 $a_\tau = \frac{\mathrm{d}v}{\mathrm{d}t}$ 反映了速度\textbf{大小（速率）}随时间的变化率。
\item $\frac{v^2}{\rho}\vec{n}$: \textbf{法向加速度矢量（Normal / Centripetal Acceleration）}。通常记作 $\vec{a}_n$。其大小 $a_n = \frac{v^2}{\rho}$ 反映了速度\textbf{方向}随时间的变化率（即向心加速度）。
\end{itemize}
\end{definition}
